% chapters/ch04_hst_n.tex
% Chapter 4: HST-N — Human State Topology for NRMO
% Source: NRMO_v5_updated.pdf, Chapter 4 (lines 5640-5866 of v5_layout.txt)
% This chapter's content overlaps with Section 2.9 (State Labels) by design;
% Chapter 4 is the standalone HST-N reference, while Section 2.9 is the
% Secretary-Console-side interface.

\chapter{HST-N: Human State Topology for NRMO}
\label{ch:hst-n}

\nrmobox{Non-Operational Charter}{
HST-N \textbf{does not generate actions}. It is a map for organising
psychological-state classifications --- not a manual for persuasion
or manipulation. Always used solely to supply input labels to the
SOP. Never used standalone.
}

% =====================================================================
\section{Role: Topography Map of Degrees of Freedom}
\label{sec:hstn-role-topography}

HST-N's role within the NRMO+Z+PP layered system is to provide
\emph{visibility} into the operator's state space. The output is a
topography of two surfaces:

\begin{itemize}
\item the \emph{irreversible boundary}, beyond which actions cannot
be undone;
\item the \emph{reversible region}, within which actions preserve
future degrees of freedom.
\end{itemize}

\noindent
This map is consumed by the SOP layer; HST-N never executes actions
based on it.

% =====================================================================
\section{Observation Inputs}
\label{sec:hstn-observation-inputs}

HST-N classifies signals from three sources:

\subsection{Subjective Observations}

\begin{itemize}
\item Mood self-report.
\item Fatigue self-report.
\item Sense of urgency, dread, or stuckness.
\end{itemize}

\subsection{Objective Observations}

\begin{itemize}
\item Sleep duration.
\item Work hours.
\item Recent commitments and capital state.
\item Relational signals (frequency, content, pattern).
\end{itemize}

\subsection{Behavioural Observations}

\begin{itemize}
\item Deviation from baseline patterns.
\item Boundary / blame-shifting (\textit{Blame Shifting}) at
relationship interfaces.
\item Out-of-band communication or unusual channel usage.
\end{itemize}

\subsection{SOP Output Vector}

\begin{table}[ht]
\centering
\small
\begin{tabularx}{\linewidth}{lXX}
\toprule
\textbf{Label} & \textbf{Domain} & \textbf{Action class} \\
\midrule
LOAD
& Judgement-load strong; reversible
& When high $\Rightarrow$ restrict to reversible (Category A/B). \\
VOL
& Volatility (stable / volatile / boundary)
& When $V_2$ $\Rightarrow$ refrain from boundary actions. \\
IRREV
& Irreversibility boundary
& Category $D$ requires procedure or SHUTDOWN. \\
\bottomrule
\end{tabularx}
\caption[HST-N output vector]{HST-N output vector consumed by SOP.
The labels do not encode SOP decisions; they are state
classifications.}
\label{tab:hstn-output-vector}
\end{table}

% =====================================================================
\section{State Labels}
\label{sec:hstn-state-labels}

\subsection{LOAD (Cognitive / Emotional Load)}
\label{sec:hstn-load}

\begin{table}[ht]
\centering
\small
\begin{tabularx}{\linewidth}{lXX}
\toprule
\textbf{Level} & \textbf{State} & \textbf{Common consequences} \\
\midrule
$L_0$
& Rested, stable, sufficient buffer
& Standard operation; full action space. \\

$L_1$
& Slight fatigue; reduced attention; confirmation friction
& Reduce category C/D actions; favour confirmation. \\

$L_2$
& High load, judgement degraded
& Restrict to A/B reversible; boundary watch. \\
\bottomrule
\end{tabularx}
\caption[HST-N LOAD levels]{LOAD level definitions and their
operational consequences.}
\label{tab:hstn-load}
\end{table}

\subsection{VOL (Volatility)}
\label{sec:hstn-vol}

\begin{table}[ht]
\centering
\small
\begin{tabularx}{\linewidth}{lXX}
\toprule
\textbf{Level} & \textbf{State} & \textbf{Common consequences} \\
\midrule
$V_0$
& Stable; predictable
& Standard operation. \\

$V_1$
& Volatile; partial unpredictability
& Slow tempo; favour A/B. \\

$V_2$
& Highly volatile; emotional outburst risk; irreversible-action
exposure
& Refrain from boundary actions; SHUTDOWN candidate. \\
\bottomrule
\end{tabularx}
\caption[HST-N VOL levels]{VOL level definitions. Higher VOL
indicates higher likelihood of breakdown in the situation.}
\label{tab:hstn-vol}
\end{table}

\subsection{IRREV (Irreversibility Proximity)}
\label{sec:hstn-irrev}

\begin{table}[ht]
\centering
\small
\begin{tabularx}{\linewidth}{lXX}
\toprule
\textbf{Level} & \textbf{State} & \textbf{Design intent} \\
\midrule
$I_0$
& Boundary distant; high reversibility
& Expand reversible degrees of freedom. \\

$I_1$
& Relationship / capital / agency boundary near
& Category $D$ requires procedure; consider SHUTDOWN. \\
\bottomrule
\end{tabularx}
\caption[HST-N IRREV levels]{IRREV level definitions. Near-boundary
states demand procedure-bound or halt-bound action selection.}
\label{tab:hstn-irrev}
\end{table}

% =====================================================================
\section{Irreversibility Types}
\label{sec:hstn-irrev-types}

The irreversible boundary is not one-dimensional. HST-N classifies
four types of irreversibility, each with its own avoidance principle.

\begin{table}[ht]
\centering
\small
\begin{tabularx}{\linewidth}{lXXX}
\toprule
\textbf{Type} & \textbf{What decreases} & \textbf{Common triggers} & \textbf{Avoidance principle} \\
\midrule
R-Trust
& Reputation, relational credibility
& Public statements, broken commitments, rumour
& Slow tempo; reversible expressions; explicit confirmation. \\

R-Resource
& Capital (time, money, attention)
& Over-commitment, ultimatum-driven binding, opportunity-cost
neglect
& Limit-based commitment; staged release; bounded loss caps. \\

R-Relationship
& Specific bilateral relationship integrity
& Ultimatums, public airing, accusatory framing
& Slow tempo; private confirmation; no broadcast. \\

R-Agency
& Personal autonomy and self-direction
& Coercion, dependency formation, role override
& Confirm own action; third-party check; explicit boundary. \\
\bottomrule
\end{tabularx}
\caption[Irreversibility types]{Four types of irreversibility used
by HST-N. Each type has its own avoidance principle, applied by
SOP based on HST-N classification.}
\label{tab:hstn-irrev-types}
\end{table}

\paragraph{Why four types matter.}
A single ``irreversibility'' scalar conflates fundamentally different
boundary classes. R-Trust and R-Relationship are highly socially
embedded; R-Resource is finitely measurable; R-Agency is partially
internal. The avoidance principles are domain-specific; the SOP
selects from category $A_{\mathrm{allowed}}$ based on which type
dominates.

% =====================================================================
\section{Constraints (Permitted / Prohibited)}
\label{sec:hstn-constraints}

\nrmobox{HST-N boundary constraints}{
\begin{itemize}
\item HST-N does not cross the boundary into evaluation
(``situation is good / bad'').
\item Strong observations are reported as classifications, never as
recommendations.
\item Output flow is one-way: HST-N $\to$ SOP. No direct action.
\end{itemize}
}

% =====================================================================
\section{HST-N SOP v1.1 --- State Monitoring \& Non-Intervention
Audit}
\label{sec:hstn-sop-v11}

\subsection*{Document status}

\begin{tabular}{ll}
\textbf{Version}: & 1.1 (Silent Failure Hardening / Multi-Window
Observations) \\
\textbf{Scope}: & Monitoring Only \\
\textbf{Execution}: & This SOP does not make decisions or execute
actions.
\end{tabular}

\subsection{Role}
\label{sec:hstn-v11-role}

\begin{itemize}
\item HST-N functions as the \emph{State Monitoring} layer.
\item Purpose: classify subject state into Low / High thresholds.
\item Capabilities: Observation, Classification, Logging.
\item Prohibitions: Judgement, Proposal, Persuasion,
Permitted-action issuance, SOP execution.
\end{itemize}

\subsection{State Definitions (Aligned with NRMO SOP)}
\label{sec:hstn-v11-state-definitions}

HST-N inherits state-transition definitions from the NRMO SOP
(asset SOP v1.2.1). State transitions are tracked for audit
purposes only:

\begin{itemize}
\item \textbf{Caution}: 4-month window (SOP-relevant).
\item \textbf{Risk-Off (Ruin)}: 12-month window (SOP-relevant).
\end{itemize}

\subsection{Negative Assertion (Audit-Only Output)}
\label{sec:hstn-v11-negative-assertion}

HST-N publishes a quarterly ``negative assertion'' record. The
pattern is structurally distinct from a positive assertion: it
states only that under the current thresholds and observation
windows, \emph{no anomaly has been detected}.

\begin{lstlisting}[basicstyle=\ttfamily\small,frame=single]
HST-N Negative Assertion (Quarterly)
- No anomaly detected under current thresholds.
- Observation windows: 4m / 12m (+ internal)
- Data availability: OK / Partial / Hold
\end{lstlisting}

\noindent
\textbf{Note}: this output is for audit only and does not influence
decision-making at the SOP layer.

\subsection{Observation Cycle (Non-Intervention)}
\label{sec:hstn-v11-observation-cycle}

The HST-N observation cycle is layered:

\subsubsection{SOP-Relevant Windows}

\begin{itemize}
\item \textbf{4-month window}: state classification feeding SOP
caution detection.
\item \textbf{12-month window}: state classification feeding SOP
risk-off detection.
\end{itemize}

\subsubsection{Audit-Only Windows}

\begin{itemize}
\item \textbf{1-month window}: short-term audit observation.
\item \textbf{6-month window}: medium-term audit observation.
\end{itemize}

\noindent
The audit-only windows do not feed the SOP directly. They exist for
post-hoc analysis of monitoring effectiveness.

\subsection{Red Flag Codes (Classification, Non-Intervention)}
\label{sec:hstn-v11-red-flag}

Red Flag codes are observational classifications, not actions. They
are forwarded to Type ZERO and SOP for downstream handling.

\begin{table}[ht]
\centering
\small
\begin{tabularx}{\linewidth}{llX}
\toprule
\textbf{Code} & \textbf{Name} & \textbf{Description} \\
\midrule
RF0 & Low-Grade Noise
& Mild signal (logged, no SOP action). \\
RF1 & Volatility Spike
& 4-month asset / benchmark volatility elevated. \\
RF2 & Halt / Data Hold
& Shutdown signal: data not reliably available. \\
RF3 & Liquidity Gap
& Bid-Ask gap, resource gap, response gap. \\
RF4 & Correlation Break
& Cross-domain signal de-coupling (observation only). \\
\bottomrule
\end{tabularx}
\caption[HST-N Red Flag codes]{HST-N Red Flag classification codes.
Codes are non-judgemental classifications passed to Type ZERO and
SOP for interpretation.}
\label{tab:hstn-v11-red-flag}
\end{table}

\paragraph{Log entry format.}

\begin{lstlisting}[basicstyle=\ttfamily\small,frame=single]
HST-N Red Flag Log
- Code:      RF1
- Window:    1m / 4m / 6m / 12m
- Asset:     (Ticker or Category)
- Timestamp: YYYY-MM-DD
- Note:      short text (no interpretation)
\end{lstlisting}

\subsection{Red Flag Frequency Summary
(Non-Intervention / Cumulative)}
\label{sec:hstn-v11-frequency-summary}

HST-N publishes a quarterly Red Flag frequency summary (audit only):

\begin{itemize}
\item \textbf{Per-code counts}: how many times each Red Flag code
fired.
\item \textbf{Distribution}: e.g., RF1 = 1, RF2 = 0, etc.
\item \textbf{Audit channel}: forwarded for review, not action.
\end{itemize}

\begin{lstlisting}[basicstyle=\ttfamily\small,frame=single]
HST-N Frequency Summary (Quarterly)
- RF0: 5
- RF1: 1
- RF2: 0
- RF3: 0
- RF4: 1
\end{lstlisting}

\subsection{Hold Mode (Data Reliability)}
\label{sec:hstn-v11-hold-mode}

When benchmark or asset data degrades, HST-N publishes a Hold
record:

\begin{itemize}
\item \texttt{OK} --- data reliable, full observation cycle proceeds.
\item \texttt{Partial} --- partial coverage, classification flagged
partial.
\item \texttt{Hold} --- data unreliable, classification suspended,
SOP notified.
\end{itemize}

\noindent
HST-N's ``Hold'' is a notification to SOP; SOP retains decision
authority.

\subsection{Update Conditions}
\label{sec:hstn-v11-update-conditions}

\begin{itemize}
\item HST-N SOP is updated only when the upstream NRMO SOP
changes.
\item Discretionary updates are prohibited.
\item All updates require explicit external re-approval.
\end{itemize}

\subsection{Prohibitions (Confirmation)}
\label{sec:hstn-v11-prohibitions}

\begin{itemize}
\item State-transition execution: \textbf{prohibited}.
\item Direct action / proposal: \textbf{prohibited}.
\item TTM / PP execution: \textbf{prohibited}.
\item Permitted-action issuance: \textbf{prohibited}.
\end{itemize}

\paragraph{Design principle.}
\textit{HST-N observes; HST-N does not decide.}

% =====================================================================
\section*{Connections to Other Parts}

\begin{itemize}
\item HST-N's output vector (LOAD / VOL / IRREV) is consumed by
the Secretary Console (Chapter~\ref{ch:secretary-modules},
specifically Section~\ref{sec:state-labels-c2}) as SOP input.
\item The Red Flag codes RF0--RF4 connect to SOP defensive
procedures defined in
Section~\ref{sec:hstn-red-flag-c2}.
\item The 4-month and 12-month windows correspond to the
Investment SOP's Caution / Risk-Off triggers
(Chapter~\ref{ch:investment-sop}).
\item In the simulator (Part~VIII), HST-N's classification function
is the cognitive analogue of the MAPLayer's mode classifier
(Chapter~\ref{ch:vnext-state}).
\item HST-N's Negative Assertion pattern is the structural
counterpart of the autonomy-preserving choice-set output of APCSO
(Section~\ref{sec:apcso-formal}): both produce non-prescriptive
outputs that defer authority to the human decision layer.
\end{itemize}
